% 贯穿图的四种表示。边编号 e1={1,2}, e2={1,3}, e3={2,3}, e4={2,4}, e5={3,4}；
% 关联矩阵对每条边取“小编号→大编号”的记录方向：起点 1，终点 -1。
\begin{tikzpicture}[>=Stealth,
  cell/.style={draw=FigureGrid,line width=.5pt,minimum width=5.6mm,minimum height=5mm,inner sep=0pt,anchor=center},
  lab/.style={text=FigureMuted,minimum width=5.6mm,minimum height=5mm,inner sep=0pt},
  one/.style={fill=FigureBlueFill},
  neg/.style={fill=FigureGreenFill},
  hd/.style={draw=FigureStroke,fill=FigureBlueFill,line width=.6pt,minimum width=5.6mm,minimum height=5mm,inner sep=0pt},
  nb/.style={draw=FigureStroke,fill=white,line width=.6pt,minimum width=5.6mm,minimum height=5mm,inner sep=0pt},
  ttl/.style={anchor=west}]
% (a) 边列表
\node[ttl] at (-.3,.75) {(a) 边列表};
\matrix[matrix of math nodes,nodes={cell},column sep=-.5pt,row sep=-.5pt,anchor=north west,
  column 1/.style={nodes={minimum width=8mm}},column 2/.style={nodes={minimum width=12mm}}] (EL) at (0,.35) {
  |[lab,minimum width=8mm]| \text{边} & |[lab,minimum width=12mm]| \text{端点} \\
  e_1 & \{1,2\} \\ e_2 & \{1,3\} \\ e_3 & \{2,3\} \\ e_4 & \{2,4\} \\ e_5 & \{3,4\} \\};
% (b) 邻接表
\begin{scope}[xshift=4.2cm]
\node[ttl] at (-.3,.75) {(b) 邻接表};
\foreach \v/\nbrs [count=\r] in {1/{2,3},2/{1,3,4},3/{1,2,4},4/{2,3}}{
  \node[hd] (h\v) at (0,{.1-(\r-1)*.75}) {$\v$};
  \foreach \u [count=\c] in \nbrs {\node[nb] (n\v\c) at ({.55+\c*.75},{.1-(\r-1)*.75}) {$\u$};}
  \draw[->,FigureStroke,line width=.6pt] (h\v.east)--(n\v1.west);
}
\end{scope}
% (c) 邻接矩阵
\node[ttl] at (-.3,-3.6) {(c) 邻接矩阵$\bm A$};
\matrix[matrix of math nodes,nodes={cell},column sep=-.5pt,row sep=-.5pt,anchor=north west] (AM) at (0,-4.0) {
  |[lab,draw=none]| & |[lab]|1 & |[lab]|2 & |[lab]|3 & |[lab]|4 \\
  |[lab]|1 & 0 & |[one]|1 & |[one]|1 & 0 \\
  |[lab]|2 & |[one]|1 & 0 & |[one]|1 & |[one]|1 \\
  |[lab]|3 & |[one]|1 & |[one]|1 & 0 & |[one]|1 \\
  |[lab]|4 & 0 & |[one]|1 & |[one]|1 & 0 \\};
% (d) 关联矩阵
\begin{scope}[xshift=4.2cm]
\node[ttl] at (-.3,-3.6) {(d) 关联矩阵$\bm B$};
\matrix[matrix of math nodes,nodes={cell,minimum width=7mm},column sep=-.5pt,row sep=-.5pt,anchor=north west] (BM) at (0,-4.0) {
  |[lab,draw=none,minimum width=5.6mm]| & |[lab,minimum width=7mm]|e_1 & |[lab,minimum width=7mm]|e_2 & |[lab,minimum width=7mm]|e_3 & |[lab,minimum width=7mm]|e_4 & |[lab,minimum width=7mm]|e_5 \\
  |[lab,minimum width=5.6mm]|1 & |[one]|1 & |[one]|1 & 0 & 0 & 0 \\
  |[lab,minimum width=5.6mm]|2 & |[neg]|-1 & 0 & |[one]|1 & |[one]|1 & 0 \\
  |[lab,minimum width=5.6mm]|3 & 0 & |[neg]|-1 & |[neg]|-1 & 0 & |[one]|1 \\
  |[lab,minimum width=5.6mm]|4 & 0 & 0 & 0 & |[neg]|-1 & |[neg]|-1 \\};
\end{scope}
\end{tikzpicture}
